\section{Related work}

\paragraph{Text-guided synthesis.}

Text-guided image synthesis has been widely studied in the context of GANs~\citep{goodfellow2014generative}. Typically, a conditional model is trained to reproduce samples from given paired image-caption datasets~\citep{zhu2019dm,tao2020df}, leveraging attention mechanisms~\citep{xu2018attngan} or cross-modal contrastive approaches~\citep{zhang2021cross,ye2021improving}. More recently, impressive visual results were achieved by leveraging large scale auto-regressive~\citep{ramesh2021zero,yu2022scaling} or diffusion models~\citep{ramesh2022hierarchical,saharia2022photorealistic,nichol2021glide,rombach2021highresolution}.

Rather than training conditional models, several approaches employ test-time optimization to explore the latent spaces of a pre-trained generator~\citep{crowson2022vqgan,murdock2021bigsleep,crowson2021diffusion}. These models typically guide the optimization to minimize a text-to-image similarity score derived from an auxiliary model such as CLIP~\citep{radford2021learning}.

Moving beyond pure image generation, a large body of work explores the use of text-based interfaces for image editing~\citep{patashnik2021styleclip,abdal2021clip2stylegan,avrahami2022blended}, generator domain adaptation~\citep{gal2021stylegan,kim2022diffusionclip}, video manipulation~\citep{tzaban2022stitch,bar2022text2live}, motion synthesis~\citep{tevet2022motionclip,petrovich22temos}, style transfer~\citep{kwon2021clipstyler,liu2022name} and even texture synthesis for 3D objects~\citep{text2mesh}.

Our approach builds on the open-ended, conditional synthesis models. Rather than training a new model from scratch, we show that we can expand a frozen model's vocabulary and introduce new pseudo-words that describe specific concepts. 

 
\paragraph{GAN inversion.}
Manipulating images with generative networks often requires one to find a corresponding latent representation of the given image, a process referred to as \textit{inversion}~\citep{zhu2016generative,xia2021gan}.
In the GAN literature, this inversion is done through either an optimization-based technique~\citep{abdal2019image2stylegan,abdal2020image2stylegan++,zhu2020improved,gu2020image} or by using an encoder~\citep{richardson2020encoding,zhu2020domain,pidhorskyi2020adversarial,tov2021designing}. Optimization methods directly optimize a latent vector, such that feeding it through the GAN will re-create a target image. Encoders leverage a large image set to train a network that maps images to their latent representations. 

In our work, we follow the optimization approach, as it can better adapt to unseen concepts. Encoders face harsher generalization requirements, and would likely need to be trained on web-scale data to offer the same freedom. We further analyze our embedding space in light of the GAN-inversion literature, outlining the core principles that remain and those that do not.

\paragraph{Diffusion-based inversion.}
In the realm of diffusion models, inversion can be performed \naive{ly} by adding noise to an image and then de-noising it through the network. However, this process tends to change the image content significantly. \citet{choi2021ilvr} improve inversion by conditioning the denoising process on noised low-pass filter data from the target image.
\citep{dhariwal2021diffusion} demonstrate that the DDIM~\citep{song2020denoising} sampling process can be inverted in a closed-form manner, extracting a latent noise map that will produce a given real image. In DALL-E 2~\citep{ramesh2022hierarchical}, they build on this method and demonstrate that it can be used to induce changes in the image, such as cross-image interpolations or semantic editing. The later relies on their use of CLIP-based codes to condition the model, and may not be applicable to other methods.

Whereas the above works invert a given \textit{image} into the model's latent space, we invert a user-provided \textit{concept}. Moreover, we represent this concept as a new pseudo-word in the model's vocabulary, allowing for more general and intuitive editing.


\paragraph{Personalization.}
Adapting models to a specific individual or object is a long-standing goal in machine learning research. Personalized models are typically found in the realms of recommendation systems~\citep{benhamdi2017personalized,fernando2018artwork,martinez2009s,cho2002personalized} or in federated learning~\citep{mansour2020three,jiang2019improving,fallah2020personalized,shamsian2021personalized}.

More recently, personalization efforts can also be found in vision and graphics. There it is typical to apply a delicate tuning of a generative model to better reconstruct specific faces or scenes~\citep{semantic2019bau,roich2021pivotal,alaluf2021hyperstyle,dinh2022hyperinverter,cao2022authentic,nitzan2022mystyle}. 

Most relevant to our work is PALAVRA~\citep{cohen2022my}, which leverages a pre-trained CLIP model for retrieval and segmentation of personalized objects. PALAVRA identifies pseudo-words in the textual embedding space of CLIP that refer to a specific object. These are then used to describe images for retrieval, or in order to segment specific objects in a scene. However, their task and losses are both discriminative, aiming to separate the object from other candidates. As we later show (\Cref{fig:baseline_comp}), their approach fails to capture the details required for plausible reconstructions or synthesis in new scenes.
